% !TEX root = ../FundationsDataScience.tex

\chapter{Sparse Regularization}
\label{chap-sparse-regul}
\chapteroverview{Many signals admit accurate approximations using only a small number of coefficients in a suitable basis or dictionary. Sparse regularization turns this observation into a reconstruction principle for noisy or incomplete measurements. We move from counting nonzero coefficients to a convex penalty, distinguish analysis and synthesis models, derive thresholding and iterative algorithms, and apply them to deconvolution and inpainting.}

The main references for this chapter are~\cite{mallat2008wavelet,starck2015sparse,scherzer2009variational}.


%%%%%%%%%%%%%%%%%%%%%%%%%%%%%%%%%%%%%%%%%%%%%%%%%%%%%%%%%%%%%%%%%%%%%%
%%%%%%%%%%%%%%%%%%%%%%%%%%%%%%%%%%%%%%%%%%%%%%%%%%%%%%%%%%%%%%%%%%%%%%
%%%%%%%%%%%%%%%%%%%%%%%%%%%%%%%%%%%%%%%%%%%%%%%%%%%%%%%%%%%%%%%%%%%%%%
\section{Sparsity Priors}

%%%%%%%%%%%%%%%%%%%%%%%%%%%%%%%%%%%%%%%%%%%%%%%%%%%%%%%%%%%%%%%%%%%%%%
\subsection{Ideal Sparsity Prior}

Chapter \ref{chap-approximation} shows that an orthonormal basis $\Bb = \{ \psi_m \}_m$ adapted to an image class $\Th$ can approximate each $f\in\Th$ using relatively few atoms.

The $\lzero$ prior measures representation complexity by counting the nonzero coefficients:
\eq{
	\Jz(f) \eqdef \#\enscond{m}{\dotp{f}{\psi_m} \neq 0}
	\qwhereq
	x_m = \dotp{f}{\psi_m}.
}
This count is also called the $\lzero$ pseudonorm:
\eq{
	\normz{x} \eqdef \Jz(f).
}
It measures sparsity directly, without accounting for the magnitudes of the nonzero coefficients.

Natural images usually have many nonzero coefficients, but can often be approximated by a function $f_M$ with a small count $M=\Jz(f_M)$. Retaining coefficients above a threshold gives a best $M$-term approximation, with $M$ determined by that threshold:
\eq{
	f_M = \sum_{ |\dotp{f}{\psi_m}|>T } \dotp{f}{\psi_m} \psi_m
	\qwhereq
	M = \#\enscond{m}{|\dotp{f}{\psi_m}|>T}.
}
For suitable wavelet bases and the bounded-variation image classes specified in Section \ref{sec-signal-models}, the error $\norm{f-f_M}$ satisfies quantitative decay estimates. These estimates motivate approximating natural images $f$ by functions with small $\Jz$.

Figure \ref{fig-sparsity-wav} displays a natural image in the wavelet basis $\psi_m = \psi_{j,n}^\om$, indexed by $m=(j,n,\om)$. Most coefficients $\dotp{f}{\psi_m}$ have small magnitude, so retaining the few large coefficients captures much of the image.

\myfigure{
\tabdeux{
\image{variational}{.3}{sparsity-image}&
\image{variational}{.3}{sparsity-wav}\\
Image $f$ & Coefficients $\dotp{f}{\psi_m}$
}
}{%
    Most wavelet coefficients of a natural image have small magnitude.
}{fig-sparsity-wav}


%%%%%%%%%%%%%%%%%%%%%%%%%%%%%%%%%%%%%%%%%%%%%%%%%%%%%%%%%%%%%%%%%%%%%%
\subsection{Convex Relaxation}

The ideal sparsity prior $\Jz$ is nonconvex, which makes optimization difficult. For example, if $f$ and $g$ have disjoint nonempty coefficient supports in $\Bb$, then $\Jz( (f+g)/2 ) = \Jz(f)+\Jz(g)$, violating convexity of $\Jz$.

In a general inverse problem, minimizing $\Jz$ entails a combinatorial search over possible coefficient supports.

To approximate the ideal prior $\Jz$, consider the $\ell^q$ family with $q>0$:
\eq{
	J_q(f) = \displaystyle \sum_m |\dotp{f}{\psi_m}|^q.
}
As shown in Figure \ref{fig-sparsity-lq}, the unit balls in $\RR^2$ become concentrated near the coordinate axes as $q\downarrow0$. For each fixed finite-dimensional vector, $J_q(f)\to\Jz(f)$; the limiting shape of these bounded unit balls should not be confused with the unbounded set $\{x:\norm{x}_0\leq1\}$. Small values of $q$ favor sparsity.


\myfigure{
\image{variational}{1}{sparsity-lq-balls}
}{%
	$\ell^q$ balls $\enscond{x}{J_q(x) \leq 1}$ for varying $q$.
}{fig-sparsity-lq}


The prior $J_q$ is convex exactly when $q \geq 1$. The smallest exponent yielding a convex prior, $q=1$, gives the $\lun$ prior $\Ju$:
\eql{\label{eq-dfn-lun-prior}
	\Ju(f) = \norm{(\dotp{f}{\psi_m})}_1 = \sum_m |\dotp{f}{\psi_m}|.
}
In the following, we consider discrete orthonormal bases $\Bb = \{\psi_m\}_{m=0}^{N-1}$ of $\RR^N$.


%%%%%%%%%%%%%%%%%%%%%%%%%%%%%%%%%%%%%%%%%%%%%%%%%%%%%%%%%%%%%%%%%%%%%%
\subsection{Sparse Regularization and Thresholding}

Given an orthonormal basis $\{\psi_m\}_m$ of $\RR^N$, regularization-based denoising \eqref{eq-regul-denoising} can be written using the sparsity priors $\Jz$ and $\Ju$ as
\eq{
	f^\star \in \uargmin{g \in \RR^N} \frac{1}{2} \norm{f-g}^2 + \la J_q(g)
}
for $q=0$ or $q=1$, writing $J_0=\Jz$. Orthonormality separates the objective into coefficientwise terms:
\eq{
	f^\star = \sum_m x^\star_m \psi_m
}
\eq{
	\qwhereq
	x^{\star} \in \uargmin{y \in \RR^N} \sum_m \frac{1}{2} |x_m-y_m|^2 + \la |y_m|^q
}
Here $x_m \eqdef \dotp{f}{\psi_m}$ and $y_m \eqdef \dotp{g}{\psi_m}$. For $q=0$, we adopt the convention
\eq{
	\foralls u \in \RR, \quad
	|u|^0=
	\choice{
		0 \qifq u=0, \\ 1 \quad \text{otherwise.} 
	}
}
Each coefficient of the denoised image solves a one-dimensional optimization problem
\eql{\label{eq-var-1D-spars}
	x_m^\star \in \uargmin{u \in \RR}  \frac{1}{2} |x_m-u|^2 + \la |u|^q
}
The following proposition gives a closed-form solution by thresholding.

\begin{prop}\label{prop-equiv-sparse-thresh}
One may choose the minimizer
\eql{\label{eq-equiv-sparse-thresh}
	x^{\star}_m = S_T^q( x_m )
	\qwhereq 
	\choice{
		T = \sqrt{2 \la} \qforq q=0,\\
		T = \la \qforq q=1,
	}
}
where 
\eql{\label{eq-hard-thresh}
	\foralls u \in \RR, \quad
	S_T^0(u) \eqdef
	\choice{
		0 \qifq |u| < T, \\
		u \quad\text{otherwise}	
	}
} 
is hard thresholding, with the alternative tie convention discussed below, and
\eql{\label{eq-soft-thresh}
	\foralls u \in \RR, \quad
	S_T^1(u) \eqdef  \sign(u) ( |u|-T )_+
} 
is the soft thresholding introduced in~\eqref{eq-soft-thresh-denoise}. For hard thresholding at $|u|=T$, both $0$ and $u$ minimize the objective; the displayed rule selects $u$. Soft thresholding gives the unique minimizer.
\end{prop}

	
\begin{proof}
For $q=0$, a zero input is minimized at zero. For a nonzero input, the best nonzero candidate is $u=x_m$, with cost $\lambda$, whereas $u=0$ has cost $|x_m|^2/2$. Comparing the two costs gives the threshold $\sqrt{2\lambda}$ and the stated tie.
For $q=1$, the optimality condition is $0\in u-x_m+\lambda\partial|u|$. If $u>0$, this gives $u=x_m-\lambda$; if $u<0$, it gives $u=x_m+\lambda$; and $u=0$ is optimal exactly when $|x_m|\leq\lambda$. These cases give soft thresholding.

\end{proof}

\begin{figure}[tbp]
\centering
{\small\sffamily\color{BookMuted}Panel 1\par}\smallskip
\BookDrawing{tikz/sparse-regularization-hard-objective}
\par\medskip
{\small\sffamily\color{BookMuted}Panel 2\par}\smallskip
\BookDrawing{tikz/sparse-regularization-soft-objective-1}
\caption{Panel 1: the objective $\norm{\cdot-y}^2+T^2 \norm{\cdot}_0$.
    Panels 2--5: the scalar objective $F(x) \eqdef \frac{1}{2}|x-y|^2+\la |x|$ for increasing $\la$.}
\label{fig-varspars}
\end{figure}
\begin{figure}[tbp]
\BookContinuedFigure
\centering
{\small\sffamily\color{BookMuted}Panel 3\par}\smallskip
\BookDrawing{tikz/sparse-regularization-soft-objective-2}
\par\medskip
{\small\sffamily\color{BookMuted}Panel 4\par}\smallskip
\BookDrawing{tikz/sparse-regularization-soft-objective-3}
\par\medskip
{\small\sffamily\color{BookMuted}Panel 5\par}\smallskip
\BookDrawing{tikz/sparse-regularization-soft-objective-4}
\caption{Panel 1: the objective $\norm{\cdot-y}^2+T^2 \norm{\cdot}_0$.
    Panels 2--5: the scalar objective $F(x) \eqdef \frac{1}{2}|x-y|^2+\la |x|$ for increasing $\la$.}
\label{fig-varspars-continued}
\end{figure}




Transforming the thresholded coefficients back to the image domain gives
\eq{
	f_{\la,q} = \sum_m S_T^q( \dotp{f}{\psi_m} ) \psi_m.
}
For Gaussian white noise $w$ of variance $\si^2$, thresholds can be chosen from the noise level; see Section \ref{sec-nl-denoising-thresh}. The universal threshold $T=\sigma\sqrt{2\log N}$ controls the largest noise coefficients and gives asymptotic risk guarantees under the signal assumptions of Section \ref{subsec-minimax-denoise}. Common empirical choices are $T \approx 3\si$ for hard thresholding ($\lzero$ regularization) and $T \approx 3\si/2$ for soft thresholding ($\lun$ regularization); see Figure \ref{fig-wavthresh-2d}.



% Bayesian interpretation:  $\dotp{f}{\psi_m}$ i.i.d. $ \sim$ generalized Gaussian


%%%%%%%%%%%%%%%%%%%%%%%%%%%%%%%%%%%%%%%%%%%%%%%%%%%%%%%%%%%%%%%%
%%%%%%%%%%%%%%%%%%%%%%%%%%%%%%%%%%%%%%%%%%%%%%%%%%%%%%%%%%%%%%%%
%%%%%%%%%%%%%%%%%%%%%%%%%%%%%%%%%%%%%%%%%%%%%%%%%%%%%%%%%%%%%%%%
\section{Sparse Regularization of Inverse Problems}
\label{sec-sparse-ip}

Using the $\lun$ prior in an orthonormal basis $\{\psi_m\}_m$ of $\RR^N$, with $\Ju$ defined in \eqref{eq-dfn-lun-prior}, gives the convex inverse problem
\eql{\label{eq-bpdn}
	f_\la \in \uargmin{f \in \RR^N} \frac{1}{2}\norm{y-\Phi f}^2 + \la \sum_m |\dotp{f}{\psi_m}|.
}
Chen, Donoho, and Saunders introduced this basis pursuit denoising formulation in \cite{chen-basis-pursuit}. Section \ref{subsec-proximal-ip} derives an iterative thresholding algorithm for~\eqref{eq-bpdn}.

%%%
\paragraph{Analysis vs. synthesis priors.}

For a nonorthogonal dictionary $\Psi=\{\psi_m\}_{m=1}^Q$, which is redundant if $Q>N$, there are two distinct extensions of~\eqref{eq-bpdn}. Each uses a convex prior $J$ in
\eql{\label{eq-regul-generic}
	f_\la \in \uargmin{f \in \RR^N} \frac{1}{2}\norm{y-\Phi f}^2 + \la J(f).
}
We also use the dictionary symbol for its synthesis operator; its adjoint is the analysis operator:
\eq{
	\Psi : x \in \RR^Q \mapsto \Psi x = \sum_m x_m \psi_m
	\qandq
	\Psi^* : f \in \RR^N \mapsto ( \dotp{f}{\psi_m} )_{m=1}^Q \in \RR^Q.
}

The analysis prior penalizes the sum of the magnitudes of correlations with the dictionary atoms:
\eql{\label{eq-sparsity-analysis}
	J_1^{\text{A}}(f) \eqdef \sum_m |\dotp{f}{\psi_m}| = \norm{\Psi^* f}_1.
}
The synthesis prior minimizes the coefficient norm over expansions of $f$ in $\Psi$:
\eql{\label{eq-sparsity-synthesis}
	J_1^{\text{S}}(f) \eqdef \umin{x \in \RR^Q, \Psi x = f} \norm{x}_1.
}
The two priors agree for an orthonormal basis, but generally differ for redundant dictionaries. Analysis regularization often requires primal--dual algorithms, as detailed in Chapter~\ref{chap-conv-duality}. Its recovery theory depends on the dictionary and on the structure of the vanishing analysis coefficients. The synthesis formulation leads directly to the coefficient-space algorithms studied below. We assume the dictionary spans $\RR^N$; otherwise the synthesis prior is $+\infty$ outside its span.

We now focus on synthesis regularization $J=J_1^{\text{S}}$, and rewrite~\eqref{eq-regul-generic} as $f_\la = \Psi x_\la$ where $x_\la$ is any solution of the following basis pursuit denoising problem
\eql{\label{eq-lasso-lagr-ip}
	x_\la \in \uargmin{x \in \RR^Q} \frac{1}{2\la} \norm{y-Ax}^2 + \norm{x}_1
} 
where the measurement matrix in coefficient space is
\eq{
	A \eqdef \Phi \Psi \in \RR^{P \times Q}. 
}
For consistent data $y\in\Im(A)$, the limit $\la\downarrow0$ leads to the constrained problem
\eql{\label{eq-lasso-constr-ip}
	x^\star \in \uargmin{A x = y} \norm{x}_1
} 
and the signal is recovered as $f^\star=\Psi x^\star \in \RR^N$.


% This can be recasted as a convex linear program, which can in turn by solved by various solver such as simplex, interior points, or Douglas-Rachford iterations.

\myfigure{
\tabtrois{
\image{invpbm}{.3}{optim-geom-l1}&
\image{invpbm}{.3}{optim-geom-l2}&
\image{invpbm}{.3}{optim-geom-l1noisy}\\
$\umin{\Phi f = y} \normu{\Psi^* f}$ &
$\umin{\Phi f = y} \norm{f}^2$ &
$\umin{\norm{\Phi f-y} \leq \epsilon} \normu{\Psi^* f}$
}
}{%
	Geometry of convex optimization problems. %
}{fig-optim-geom}


 
%%%%%%%%%%%%%%%%%%%%%%%%%%%%%%%%%%%%%%%%%%%%%%%%%%%%%%%%%%%%%%%%
%%%%%%%%%%%%%%%%%%%%%%%%%%%%%%%%%%%%%%%%%%%%%%%%%%%%%%%%%%%%%%%%
%%%%%%%%%%%%%%%%%%%%%%%%%%%%%%%%%%%%%%%%%%%%%%%%%%%%%%%%%%%%%%%%
\section{Iterative Soft Thresholding Algorithm}
\label{subsec-proximal-ip}

We now derive iterative methods for~\eqref{eq-lasso-lagr-ip}, beginning with a comparison to its constrained counterpart.

%%%%%%%%%%%%%%%%%%%%%%%%%%%%%%%%%%%%%%%%%%%%%%%%%%%%%%%%%%%%%%%%
\subsection{Noiseless Recovery as a Linear Program}
\label{eq-linearprog-lasso}

Before treating $\la>0$, note that~\eqref{eq-lasso-constr-ip} can be written as a linear program. Split $x=x_+-x_-$ with $(x_+,x_-)\in(\RR_+^Q)^2$ and solve
\eql{\label{eq-lin-prog-lasso}
	(x_+^\star,x_-^\star) \in \uargmin{(x_+,x_-) \in (\RR_+^Q)^2} \enscond{ \dotp{x_+}{\ones_Q}+\dotp{x_-}{\ones_Q} }{ y = A(x_+ - x_-) }.
}
Then recover $x^\star=x_+^\star-x_-^\star$. For small or moderate $Q$, the linear program can be solved by simplex or interior-point methods.
%
For large imaging and learning problems, first-order methods are often preferable. One option is Douglas--Rachford splitting, developed in Section~\ref{sec-dr}.
%
The penalized problem~\eqref{eq-lasso-lagr-ip} admits a particularly simple splitting algorithm because its fidelity term is smooth. The relative computational cost of the two formulations depends on the solver and on $A$.


%%%%%%%%%%%%%%%%%%%%%%%%%%%%%%%%%%%%%%%%%%%%%%%%%%%%%%%%%%%%%%%%
\subsection{Projected Gradient Descent for \texorpdfstring{$\ell^1$}{l1}}

We first solve~\eqref{eq-lasso-lagr-ip} by projected gradient descent, which is analyzed in detail in Section~\ref{sec-proj-grad}.
%
As in~\eqref{eq-lin-prog-lasso}, we rewrite~\eqref{eq-lasso-lagr-ip} as a constrained minimization problem of the form~\eqref{eq-constr}, with the nonnegative constraint set $\Cc$ and
\eq{
	u=(u_+,u_-) \in (\RR^Q)^2, \quad
	\Cc=(\RR_+^Q)^2, 
	\qandq
	\Ee(u) = \frac{1}{2}\norm{ A(u_+-u_-) - y }^2 + \la \dotp{u_+}{\ones_Q} +\la \dotp{u_-}{\ones_Q}.
}
Projection onto $\Cc$ is easy to compute
\eq{
	\Proj_{(\RR_+^Q)^2}(u_+,u_-) = ((u_+)_\oplus,(u_-)_\oplus)
	\qwhereq
	(r)_\oplus \eqdef \max(r,0), 
}
and the gradient reads
\eq{
	\nabla \Ee(u_+,u_-) =  ( \eta + \la \ones_Q , -\eta + \la \ones_Q) 
	\qwhereq
	\eta = A^*( A(u_+-u_-) - y )
}

Denoting $\it{u} = (\it{u_+},\it{u_-})$ and $\it{x} \eqdef \it{u_+}-\it{u_-}$, the iterates of projected gradient descent~\eqref{eq-proj-grad-desc} are
\eq{
	\iit{u_+} \eqdef \pa{ \it{u_+} - \tau_\ell ( \it{\eta} + \la  ) }_\oplus
	\qandq
	\iit{u_-} \eqdef \pa{ \it{u_-} - \tau_\ell ( -\it{\eta} + \la  ) }_\oplus
} 
where $\it{\eta}\eqdef A^*(A\it{x}-y)$, and adding the scalar $\lambda$ means adding $\lambda\ones_Q$ componentwise.

Theorem~\ref{thm-proj-grad} ensures that $\it{u} \rightarrow u^\star$, a solution of~\eqref{eq-constr}, if
\eq{
	\foralls \ell, \quad
	0 < \tau_{\min} <  \tau_\ell < \tau_{\max} < \frac{1}{\norm{A}^2},
}
The smooth objective in $(u_+,u_-)$ has gradient Lipschitz constant $2\norm{A}^2$, which explains the step-size bound. Consequently, $\it{x} \rightarrow x^\star = u_+^\star-u_-^\star$, a solution of~\eqref{eq-lasso-lagr-ip}.


%%%%%%%%%%%%%%%%%%%%%%%%%%%%%%%%%%%%%%%%%%%%%%%%%%%%%%%%%%%%%%%%
\subsection{Iterative Soft Thresholding and Forward Backward}
\label{sec-ista}

Splitting positive and negative parts requires storing $2Q$ coefficients. The iterative soft thresholding algorithm (ISTA) avoids this duplication while retaining comparable convergence guarantees.
%
For a fixed step $0<\tau<2/\norm{A}^2$, ISTA converges to a minimizer. The surrogate derivation below uses the stricter bound $\tau\leq1/\norm{A}^2$, which gives a direct proof of energy decrease.

ISTA was derived by several authors~\cite{figueiredo-nowak-em,daubechies-iterated}. It is a special case of forward--backward proximal splitting~\cite{combettes-proximal}, studied in Section~\ref{sec-fb}.

We derive ISTA for the $\ell^1$ penalty by minimizing a sequence of surrogate objectives.
%
Rescale the objective in~\eqref{eq-lasso-lagr-ip} as
\eq{
	\Ee(x) \eqdef \frac{1}{2} \norm{y-Ax}^2 + \la \norm{x}_1
}
For a fixed reference point $x'$, define
\eq{
	\Ee_\tau(x,x') \eqdef \Ee(x) - \frac{1}{2}\norm{Ax-Ax'}^2 + \frac{1}{2\tau}\norm{x-x'}^2.
}
We have $\Ee_\tau(x,x)=\Ee(x)$, and the difference between the two objectives is
\eq{
	K(x,x') \eqdef - \frac{1}{2}\norm{Ax-Ax'}^2 + \frac{1}{2\tau}\norm{x-x'}^2 =
	\frac{1}{2}\dotp{ \pa{\frac{1}{\tau}\Id_Q-A^*A} (x-x') }{x-x'}.
}
The correction $K(x,x')$ is nonnegative when $\la_{\max}(A^*A) \leq 1/\tau$, where the left side is the largest eigenvalue. Equivalently, $\tau \leq 1/\norm{A}_{\text{op}}^2$, with $\norm{A}_{\text{op}} = \si_{\max}(A)$ denoting the operator norm.
%
Under this condition, $\Ee_\tau(x,x')$ is a surrogate that bounds the objective from above and agrees with it at the reference point:
\eq{
	\Ee(x) \leq \Ee_\tau(x,x'), \quad
	\Ee_\tau(x',x')=\Ee(x'), \qandq
	\Ee(\cdot)-\Ee_\tau(\cdot,x') \text{ is smooth.} 
}
Minimizing this surrogate at each iteration gives
\eql{\label{eq-ista-surrog}
	\iit{x} \eqdef \uargmin{x} \Ee_{\tau_\ell}(x,\it{x})
}
The surrogate upper bound and equality at the current iterate imply
\eq{
	\Ee(\iit{x}) \leq \Ee(\it{x}).
}

\begin{prop}
	The iterates $\it{x}$ defined by~\eqref{eq-ista-surrog} satisfy
	\eql{\label{eq-ista}
		\iit{x} = \Ss^1_{\la\tau_\ell}\pa{
			\it{x} - \tau_\ell A^*( A \it{x} - y ) 
		}
	}
	where $\Ss^1_\la(x) = (  S^1_\la(x_m) )_m$ and $S^1_\la(r) = \sign(r) (|r|-\la)_\oplus$ is the soft thresholding operator defined in~\eqref{eq-soft-thresh}.
\end{prop}
\begin{proof}
    Expanding the quadratic terms and collecting constants independent of the optimization variable gives
	\begin{align*}
		\Ee_\tau(x,x') &= \frac{1}{2}\norm{Ax-y}^2 - \frac{1}{2}\norm{Ax-Ax'}^2 + \frac{1}{2\tau}\norm{x-x'}^2 + \la \norm{x}_1 \\
			& = C + \frac{1}{2}\norm{Ax}^2-\frac{1}{2}\norm{Ax}^2+\frac{1}{2\tau}\norm{x}^2
			- \dotp{Ax}{y} + \dotp{Ax}{Ax'} - \frac{1}{\tau}\dotp{x}{x'}
			+ \la \norm{x}_1 \\
			& = C + \frac{1}{2\tau}\norm{x}^2 + \dotp{x}{ -A^*y+A^*Ax'-\frac{1}{\tau}x' } + \la \norm{x}_1 \\
			&= C' + \frac{1}{\tau} \pa{
				\frac12\norm{ x - \pa{ x'-\tau A^*(Ax'-y) } }^2 + \tau \la \norm{x}_1
			}
	\end{align*}
    Proposition~\ref{prop-equiv-sparse-thresh} therefore identifies the minimizer of $\Ee_\tau(x,x')$ as
    $\Ss^1_{\la\tau}( x' - \tau A^*( A x' - y ) )$.
\end{proof}

The iterations~\eqref{eq-ista} coincide with the forward--backward iterates~\eqref{eq-fb-split}. For the coefficient-space objective~\eqref{eq-lasso-lagr-ip}, rescaled by $\lambda$, use the splitting
\eql{\label{eq-bpdn-func}
	\Ff=\frac{1}{2} \norm{A \cdot - y}^2
	\qandq 
	\Gg=\la \norm{\cdot}_1. 
}
%
In the case~\eqref{eq-bpdn-func}, Proposition~\ref{prop-equiv-sparse-thresh} shows that $\Prox_{\rho\Gg}$ is soft thresholding at level $\rho\lambda$.



%%%%%%%%%%%%%%%%%%%%%%%%%%%%%%%%%%%%%%%%%%%%%%%%%%%%%%%%%%%%%%%%
%%%%%%%%%%%%%%%%%%%%%%%%%%%%%%%%%%%%%%%%%%%%%%%%%%%%%%%%%%%%%%%%
%%%%%%%%%%%%%%%%%%%%%%%%%%%%%%%%%%%%%%%%%%%%%%%%%%%%%%%%%%%%%%%%
\section{Example: Sparse Deconvolution}


%%%%%%%%%%%%%%%%%%%%%%%%%%%%%%%%%%%%%%%%%%%%%%%%%%%%%%%%%%%%%%%%
\subsection{Sparse Spikes Deconvolution}

Sparse spike deconvolution uses sparsity in the spatial domain, which corresponds to the orthonormal basis of Diracs $\psi_m[n] = \de[n-m]$. In seismic imaging, a simple sparse reflectivity model represents $f_0$ as a few impulses associated with changes in acoustic impedance.

In a linearized one-dimensional model that neglects multiple reflections, the subsurface reflectivity $f_0$ produces observations $y=h \star f_0 + w$. Here $h$ is the transmitted pulse, called a seismic wavelet. This use of the word differs from the orthogonal wavelet bases constructed in Chapter \ref{chap-wavelets}, although the terminology originated in seismic imaging.

The seismic wavelet $h$ is typically band-pass, balancing spatial and frequency concentration. Figure~\ref{fig-spikes-lun} shows a second derivative of a Gaussian and its Fourier transform. The narrow transmitted band illustrates how acquisition suppresses information about $f_0$.

% \myfigure{
% \image{invpbm}{.6}{seismic-data}
% }{%%
%	2D+t seismic data. %	
% }{fig-seismic-data}

Using $\lun$ regularization in the Dirac basis gives
\eq{
	 f^\star = \uargmin{f \in \RR^N} \frac{1}{2} \norm{f \star h - y}^2 + \la \sum_m |f_m|.
}
Figure \ref{fig-spikes-lun} shows the result of $\lun$ minimization with an oracle choice of $\la$ that minimizes the error $\norm{f^\star-f_0}$.


\myfigure{
\tabdeux{
\image{invpbm}{.4}{spikes-filter}&
\image{invpbm}{.4}{spikes-filter-fourier}\\
$h$ & $\hat h$ \\
% \image{invpbm}{.3}{spikes-filter-fourier-inv}&
\image{invpbm}{.4}{spikes-signal}&
\image{invpbm}{.4}{spikes-observations}\\
$f_0$ & $y=h \star f_0 + w$ \\
\image{invpbm}{.4}{spikes-pinv}&
\image{invpbm}{.4}{spikes-l1}\\
$f^+$ & $f^\star$ 
}
}{%
    Sparse spike recovery by the pseudoinverse and by $\lun$ regularization. %
}{fig-spikes-lun}

ISTA for sparse spike recovery alternates the updates
\eq{
	\tilde f^{(k)} = f^{(k)} - \tau h^\vee \star ( h \star f^{(k)}-y )
}
and
\eq{
	f^{(k+1)}_m = S_{\la \tau}^{1}(\tilde f^{(k)}_m )
}
where $h^\vee[n]=\overline{h[-n]}$ is the adjoint convolution kernel. It equals $h$ for a real symmetric filter. The step size must satisfy
\eq{
	0<\tau < 2/\norm{\Phi^* \Phi} = 2 / \umax{\om} |\hat h(\om)|^2
}
to guarantee convergence. Figure \ref{fig-spikes-ista-decay} shows convergence of both the objective values and the iterates. Objective convergence alone would not ensure convergence of the iterates when minimizers are nonunique; the forward--backward convergence theorem provides the latter guarantee here.


\myfigure{
\tabdeux{
\image{invpbm}{.4}{spikes-ista-decay-conv}&
\image{invpbm}{.4}{spikes-ista-decay-energy} \\
%\image{invpbm}{.3}{spikes-snr-lambda}\\
$\log_{10}( E(f^{(k)})/E(f^\star)-1 )$ &
$\log_{10}( \norm{f^{(k)}-f^\star}/\norm{f_0} )$
}
}{%
	Convergence of the energy and iterates under iterative soft thresholding. %
}{fig-spikes-ista-decay}


%%%%%%%%%%%%%%%%%%%%%%%%%%%%%%%%%%%%%%%%%%%%%%%%%%%%%%%%%%%%%%%%
\subsection{Sparse Wavelets Deconvolution}


Camera images can be blurred by defocus, motion during exposure, or diffraction. Assuming spatial invariance, we model the blur operator $\Phi$ by convolution:
\eq{
	y = f_0 \star h + w.
}
We use a Gaussian filter $h$ of width $\mu > 0$. On a fixed $d$-dimensional domain, the number of effectively transmitted frequencies scales roughly as $\mu^{-d}$, with a threshold-dependent factor, because higher frequencies are strongly attenuated.

Figures~\ref{fig-deconvolution-1d} and~\ref{fig-deconvolution-2d} show examples of signal and image acquisition with Gaussian blur.

Sobolev regularization~\eqref{eq-quad-regul-convol} suppresses high-frequency noise more strongly than the squared-norm penalty~\eqref{eq-regul-normsq}. Both can blur sharp transitions. TV regularization and sparsity in a wavelet basis are better adapted to signals with edges.

Figure \ref{fig-deconvolution-1d} compares wavelet and Sobolev regularization for a one-dimensional signal; Figure \ref{fig-deconvolution-2d} shows the corresponding image-deblurring results. Replacing an orthogonal wavelet basis in \eqref{eq-ista} by a translation-invariant tight frame~\eqref{eq-ti-wavframe} can reduce artifacts. For a redundant frame, synthesis ISTA acts in coefficient space with $A=\Phi\Psi$. Simply analyzing, thresholding, and synthesizing is generally not the proximal map of the analysis penalty.

\myfigure{
\tabdeux{
%\image{invpbm}{.3}{deconv1d-filter-fourier}&
\image{invpbm}{.4}{deconv1d-signal}&
\image{invpbm}{.4}{deconv1d-filter} \\
Signal $f_0$ & Filter $h$ \\
\image{invpbm}{.4}{deconv1d-observations}&
\image{invpbm}{.4}{deconv1d-l1}\\
Observation $y=h\star f_0+w$ & $\lun$ recovery $f^\star$
}
}{%
	One-dimensional deconvolution with sparsity in an orthogonal wavelet basis. %
}{fig-deconvolution-1d}

% \image{invpbm}{.3}{deconv1d-snr}

%%


\myfigure{
\tabtrois{
\image{invpbm}{.3}{deconv2d-image}&
\image{invpbm}{.3}{deconv2d-observations}&
\image{invpbm}{.3}{deconv2d-l2}\\
Image $f_0$ & Observations $y=h\star f_0+w$ & $\ldeux$ regularization \\
\image{invpbm}{.3}{deconv2d-sob}&
\image{invpbm}{.3}{deconv2d-l1}&
\image{invpbm}{.3}{deconv2d-l1inv}\\
Sobolev regularization & $\lun$ regularization & $\lun$ invariant \\

}
}{%
	Image deconvolution. %	
}{fig-deconvolution-2d}

Figure \ref{fig-deconvolution-2d-snr} plots the SNR against $\la$. Computing the SNR requires the clean reference image $f_0$; the reported experiments use the maximizing value of $\la$.

\myfigure{
\tabtrois{
\image{invpbm}{.3}{deconv2d-l2-snr}&
\image{invpbm}{.3}{deconv2d-sob-snr}&
\image{invpbm}{.3}{deconv2d-l1-snr}\\
$\ldeux$ regularization  &
Sobolev regularization  &
$\lun$ regularization  
}
}{%
	SNR as a function of $\la$. %	
}{fig-deconvolution-2d-snr}


%\matlab{matlab/inverse-deconv.m
%}{
%Deconvolution with $\lun$ regularization in a wavelet basis. The filter $\phi$ is given in \matvar{phi}, the observations in \matvar{y}, the regularization parameter $\la$ is \matvar{lambda}. The solution is given in \matvar{fspars}.
%}{inverse-deconv}



%%%%%%%%%%%%%%%%%%%%%%%%%%%%%%%%%%%%%%
\subsection{Sparse Inpainting}

This section continues the discussion in Section~\ref{sec-inpainting-variational}.

For noiseless inpainting, a small $\lambda>0$ approximates constrained $\ell^1$ minimization. Exact interpolation is obtained in the limit under the assumptions of Proposition~\ref{prop-gamma-conv-regul}.
For an orthonormal basis, the iterative thresholding algorithm~\eqref{eq-ista} takes the following form for $\tau=1$:
\eq{
	f^{(k+1)} = \sum_m S_{\la}^{1}( \dotp{ P_y(f^{(k)}) }{\psi_m} )  \psi_m
}
Figure \ref{fig-inpainting-lena} compares Sobolev inpainting with wavelet sparsity priors, including a translation-invariant frame.

\myfigure{
\tabtrois{
\image{invpbm}{.3}{inpainting-lena-image}&
\image{invpbm}{.3}{inpainting-lena-observations}& \\
Image $f_0$ & Observation $y=\Phi f_0$ & \\
\image{invpbm}{.3}{inpainting-lena-sob}&
\image{invpbm}{.3}{inpainting-lena-wavortho}&
\image{invpbm}{.3}{inpainting-lena-wavti}\\
Sobolev $f^\star$ & Ortho. wav $f^\star$ & TI. wav $f^\star$ \\

}
}{%
    Inpainting with Sobolev and wavelet sparsity priors. %
}{fig-inpainting-lena}

